\section{Experimental Details and Scope}
\label{app:details}

\subsection{Sampling and Feature Discovery}
\label{app:sampling}

Teacher sampling uses temperature 0.7, top-$p$ 0.95, and a 512-token generation cap.
The initial pool contains 16 candidates for each of 881 problems (14,096 traces).
Shared TopK SAEs are fit on pooled residual activations from raw candidates regardless of correctness; only answer-correct traces enter short--long feature comparisons.
Question-level splits contain 617 training, 132 development, and 132 test problems.
The SAE grid covers layers 10, 17, and 23; TopK sparsities 32 and 64; expansion factor 8; and 1,500 training steps, using seed 17.
Registered screening of the primary layer-17, TopK-64 dictionary yields eight features per length group.

Natural-length comparisons require at least four distinct, answer-correct, non-capped candidates per problem.
Short and long representatives come from the bottom and top 20\% length bands; medium representatives are near the median.
All conditions retain the same 881 problems.
Medium supervision achieves 66.33\% seed-mean accuracy, between short (66.67\%) and long (64.20\%).
Holm correction covers three registered natural-length comparisons.
Final-answer correctness does not certify intermediate steps.

\subsection{Intervention Screening}
\label{app:screening}

We screen short-feature enhancement and joint short-enhancement/long-direction subtraction at $\rho\in\{0.05,0.15,0.30\}$.
Each has a control matched in decoder-column count and actual perturbation norm, excluding target features from random sets.
Including unmodified generation gives 13 conditions.
The joint variant uses a fixed composite direction, not activation-gated suppression.
Interventions act at the layer-17 post-block residual position predicting each next token, beginning with the first generated token.
Conditions share prompts and per-question random uniforms but maintain independent attention caches.

Selection uses 64 development problems.
The shortest eligible condition must shorten mean length by at least 5\%, lose no more than five accuracy points, have a paired 95\% bootstrap length-difference interval below zero, and be shorter on average than its matched random control.
A subsequent 64-problem check gives 246.44 versus 209.23 mean tokens, with unmodified and selected generation both correct on 64/64 problems.
Both cohorts were previously used in SAE analysis; the accuracy threshold is not a non-inferiority test.
The full teacher pool contains 10,572 outputs; its bootstrap intervals resample problems while retaining all four candidates together.

\clearpage
\subsection{Student Training and Budget Accounting}
\label{app:training}

All conditions start from the same pretrained student.
LoRA uses rank 4, scaling factor 16, dropout 0.05, and all linear modules.
Completion-only SFT runs for one epoch over each constructed dataset, with batch size 4, gradient accumulation 1, learning rate $2\times10^{-5}$, warmup ratio 0.03, BF16, maximum sequence length of 2,048 tokens without packing.
Student seeds are 17, 42, and 73; these vary SFT, not SAE discovery.
Evaluation uses \texttt{test[50:1319]}; the first 50 test problems are reserved for smoke checks.
Final evaluation accuracy is not used to select interventions.

After deduplication and shortest-correct selection, unmodified, SAE, and random generation cover 881, 880, and 879 problems, respectively.
Intersecting these with natural-short support leaves 878 problems.
Equal-example datasets contain 878 records and require 220 optimizer steps.
Repeating complete solutions brings target-token totals to 215,643--215,722 (Table~\ref{tab:training_budgets}).
Random-control data, weights, and predictions are identical across budgets for each seed, not independent replications.

\begin{table}[htbp]
\centering
\small
\begin{tabular}{lrr}
\toprule
Supervision & Equal-example target tokens & Equal-token steps \\
\midrule
Unmodified & 188,540 & 252 \\
SAE short enhancement & 165,530 & 286 \\
Matched random & 215,643 & 220 \\
Natural short & 188,420 & 251 \\
\bottomrule
\end{tabular}
\caption{Training-budget accounting. Every equal-example condition uses 220 steps; equal-token conditions have approximately matched token totals but different step counts.}
\label{tab:training_budgets}
\end{table}

\subsection{Uncompleted Extension and Interpretation}
\label{app:scope}

A separately registered extension compares one, two, and four distinct correct responses per problem under SAE and unmodified generation on the same 878 problems.
Repeating the shortest response two or four times controls record count, question exposure, and optimizer steps, but not target-token counts.
Distinct text need not represent distinct mathematical solutions.
No complete audited results are available, so no conclusions are drawn from partial runs.

Lexical and answer-format explanations remain possible: random directions control feature count and perturbation magnitude, not semantics or activation prevalence.
End-to-end savings would also need to include SAE fitting, candidate generation, selection, and intervention overhead.
The present evidence does not establish that long traces are generally redundant.

\section{SAE Diagnostics and Proposed Analysis Protocol}
\label{app:sae_protocol}

\subsection{Existing Dictionary-Quality Measurements}
\label{app:sae_quality}

Table~\ref{tab:sae_quality} reproduces the completed SAE audit's held-out reconstruction and activation-coverage measurements.
The primary layer-17, TopK-64 SAE has explained variance 0.776 and approximately 0.18\% of features inactive on the measured test support.
The much higher inactive fractions at TopK 32 motivate inspecting training-set activity and learning curves before interpreting differences in the number of discovered features.
Test inactivity is not equivalent to a feature never activating during training.
These measurements neither establish semantic purity nor demonstrate that reconstruction quality predicts steering utility.

\begin{table}[htbp]
\centering
\small
\begin{tabular}{rrrr}
\toprule
Layer & TopK & Test explained variance & Test-inactive features (\%) \\
\midrule
10 & 32 & 0.812 & 53.52 \\
10 & 64 & 0.837 & 1.38 \\
17 & 32 & 0.747 & 36.18 \\
17 & 64 & 0.776 & 0.18 \\
23 & 32 & 0.759 & 18.16 \\
23 & 64 & 0.789 & 0.10 \\
\bottomrule
\end{tabular}
\caption{Existing measurements from the six completed SAEs, all trained with seed 17. Inactivity is measured on the audit's test support, not asserted over the full training distribution.}
\label{tab:sae_quality}
\end{table}

\subsection{Region-Specific Activation and Counterexamples (Proposed)}
\label{app:region_protocol}

The reported eight short-associated features are F7339, F10807, F17885, F27038, F22686, F11240, F24148, and F26001 in the primary dictionary.
All have reported high-activation examples at \texttt{Answer}; their IDs are dictionary-specific.
For each feature, a region-specific analysis would report activation mass in the answer marker and suffix, body-token means, marker occurrence counts, and activation per marker occurrence.
Paired effect estimates before and after excluding these regions would be displayed with uncertainty and support counts.
Global sums would be diagnostic only, because they also depend on trace length.
Discovery-set nuisance adjustments would be frozen before application to separate questions.

Meaning-preserving rewrites would be checked for the same numerical dependencies and answer, while format-preserving incorrect variants would change a controlled calculation.
Non-reasoning examples with injected marker tokens would test whether arithmetic content is necessary for activation.
Results would include both successful and failed counterexamples rather than only the most interpretable feature examples.
Deleting a word changes subsequent positions and states; aligned-span comparisons and matched control edits would therefore accompany lexical scrubbing.

For early-prefix prediction, models using lexical counts and absolute position would serve as controls for the SAE representation.
Dense residual features and PCA features would be additional representation baselines under the same question-grouped splits and development tuning budget.
Only traces still generating at a fixed landmark contribute to remaining-length prediction, with early-termination coverage reported separately.
Prediction quality measures readable information; it does not establish that the predicted feature causes the outcome.

\subsection{Readback and Geometric Controls (Proposed)}
\label{app:readback_protocol}

Let $\mathcal{F}$ denote the selected feature set and $\Delta z=z'-z$ the same-state readback difference.
We would report signed changes in the selected features, $\|\Delta z_{\mathcal F}\|_2$, and $\|\Delta z_{\overline{\mathcal F}}\|_2$ separately.
An off-target-to-target ratio would be accompanied by its component magnitudes and the fraction of states with a near-zero target change, avoiding interpretation driven by an unstable denominator.
Because aggregate off-target norms depend on how many coordinates are included, per-coordinate distributions and matched-random controls would also be reported.
The KL convention would be fixed as $D_{\mathrm{KL}}(p_0\|p_\delta)$ using the actual next-token distributions after all remaining layers.

Decoder cosine similarities, activation correlations, and singular values would characterize redundancy within the selected set.
A discovered answer-format subspace could further decompose the composite direction $v$ into its projection and residual:
\begin{equation}
  v_{\mathrm{format}}=P_{\mathrm{format}}v,
  \qquad v_{\mathrm{residual}}=(I-P_{\mathrm{format}})v.
\end{equation}
The subspace would be estimated on discovery data and tested on different questions.
Interventions would disclose both component norms and any normalization multiplier; a negligible residual would not be amplified into a nominally equal-norm control.
Persistence of an effect in the residual would test geometric separability from this estimated format subspace, not certify a pure reasoning mechanism.

\subsection{Continuation, Robustness, and Student Checks (Proposed)}
\label{app:continuation_protocol}

Continuation tests would clone caches from the same naturally sampled 32- or 64-token prefix and record prefixes that end too early or already contain an answer marker.
Shared random uniforms reduce sampling variation but do not force identical paths after conditions diverge.
Windowed or single-position interventions would evaluate when an effect arises; rules triggered by future trace length would not be used for online steering.
Sign reversal need not produce a symmetric response, so dose curves, correctness, and coverage would be reported together.

Additional SAE seeds 42 and 73 at the primary setting would reuse a fixed corpus and activation-sampling budget.
Feature matching would combine decoder geometry with activation profiles on separate questions and allow one-to-many correspondences.
A question- or trace-balanced token-sampling sensitivity check would test whether longer traces influence dictionary fitting primarily through their greater token count.
Reconstruction error would also be stratified by length, correctness, and position.
None of these proposed retrainings is included in the current three-seed student result.

Method-blinded paired annotation would count preserved mathematical dependencies, repeated calculations, explanations, verification, and answer formatting.
Student likelihood comparisons would separate body, calculation, and answer tokens and report segment lengths.
The final test would remain SFT on matched-support datasets and paired evaluation, rather than treating likelihood or text similarity as teaching-utility labels.
Suggested visual summaries are region-by-feature activation matrices, cleaned versus original paired effects, target/off-target dose curves, and student accuracy--length plots with uncertainty.
These figures require the corresponding completed measurements; no synthetic results are supplied for them in this draft.
